% COURSE-ASSISTANT SOURCE — FALL 2025 ARCHIVE
% Archived practice material only. This file is not evidence about Fall 2026
% assignments, deadlines, grading, submission rules, or examination coverage.
% Historical administrative instructions have been omitted. No solutions are
% included. External textbook and article figures are not included here.
\documentclass[12pt,a4paper]{article}

\input{EC2211_PSpreamble.txt}
\title{Problem Set 2}        % Sets the document title.


\begin{document}
\pagenumbering{gobble} % Stops counting the pages from this point until changed again.
\maketitle

\section{Technology transfer in the Solow model (1 p)}
(Exercise 2 in Jones chapter 5)
One explanation for China's rapid economic growth during the past several decades is its expansion of policies that encourage "technology transfer". By this, we mean policies -- such as opening up to international trade and attracting multinational corporations through various incentives -- that encourage the use and adoption in China of new ideas and new technologies. This question asks you to use the Solow model to study this scenario.

Suppose China begins in steady state. To keep the problem simple, let's assume that the sole result of these technology transfer policies is to increase $A$ by a large and permanent amount, one time. Answer the following questions:
\begin{itemize}
    \item[a.] Analyze this change using the Solow diagram. What happens to the economy over time?
    \item[b.] Draw a graph showing what happens to output in China over time. What happens to output per person in China in the long run?
    \item[c.] Draw a graph showing what happens to the growth rate of output in China over time. Explain.
    \item[d.] Discuss in a couple of sentences what your results imply about the effect of technology transfer on economic growth.
\end{itemize}

\section{A natural disaster in the Malthusian model (0.5 p)}
Consider the Malthusian model covered in lecture notes 3. Suppose that the economy is in steady state. Then, as a result of a natural disaster, a large fraction of the economy's land area becomes permanently non-arable.
\begin{itemize}
    \item[a.] Make a graph corresponding to Figure 3 in Steinsson (2025a). What happens on impact in the graph as a consequence of the natural disaster?
    \item[b.] Use the model to analyze how the population and income per capita develop over time. Present your analysis and conclusions in graphs corresponding to Figures 8 and 9 in Steinsson.
\end{itemize}


\section{A decline in research productivity (1 p)}
(Exercise 6 in Jones chapter 6)
Consider the basic Romer model presented in this chapter. Suppose the economy is on a balanced growth path and then, in the year 2030, research productivity $\zeta$ declines immediately and permanently by 10\%.
\begin{itemize}
    \item[a.] What happens to the long-run growth rate of knowledge and $y_t$?
    \item[b.] Make a graph of $y_t$ over time using a ratio scale. Explain what happens to the growth rate of $y_t$ over time.
    \item[c.] Why might reasearch productivity decline in an economy?
\end{itemize}



\section*{Recommended extra exercises }
Do not hand in solutions to the exercises below! Exercises marked by a * are solved in the textbook. Other exercises may be discussed in your sessions if time allows. 

\begin{itemize}
    \item Jones chapter 6: Exercises 1, *2 and 8
\end{itemize}

\end{document}
